\section{Generative Search Engine}
\label{sec:generative_search_engine}
This section introduces GSE (Generative Search Engine) and its system model.
GSEs are information retrieval systems that integrate web search results from queries transformed from user questions and generate synthesized answers while citing web content using large language models (LLMs)~\cite{geo-apmr+24}.
GSEs receive a question from users, perform web searches, and generate answers with LLMs using the search results.

Aggarwal et al.\ formalize GSEs and provide a system model~\cite{geo-apmr+24}.
A GSE is formalized as a function \( f_{GSE}^{model} \) that takes a user question \( q_u \) and personalization information \( P_U \) as inputs and generates a textual answer \( r \):
$ f_{GSE}^{model} := (q_u, P_U) \rightarrow r $.

GSEs consist of two components: content retrieval and answer generation.
Content retrieval collects the information necessary to generate answers from the web.
The user query \( q_u \) is converted by LLMs into multiple queries \( Q^{\prime} = \{q_1, q_2, \cdots, q_n\} \) for the web search.
These queries are sent to a search engine \( SE \), and GSEs obtain a set of web sources \( S = \{s_1, s_2, \cdots, s_m\} \) from the search results.
Each result set in $S$ for a web search query \( q_i \) typically consists of the top \( k \) web sources ranked by the search engine's metrics~\cite{poisonedrag-arxiv24, modern-information-retrieval-a-brief-overview-3320, geo-apmr+24}.


Answer generation cites the web sources obtained in the content retrieval phase and generates answers to user queries.
Web sources \( S \) are converted into a summary set \( Sum = \{Sum_1, Sum_2, \cdots, Sum_m\} \) that extracts and summarizes the content of web sources for generating answers to the user query \( q_u \).
LLMs then generate the final answer \( r \) from the summary set \( Sum \) while citing web sources \( S \).
The answer \( r \) consists of \( k \) sentences \( \{l_1, l_2, \cdots, l_k\} \), and each sentence \( l_i \) is associated with a citation set \( C_i \subseteq S, C_i = \{c_1, c_2, \cdots, c_l\} \).
Although a GSE ideally has one or more citations for each sentence, there are cases where \( C_i = \emptyset \), meaning there are no citations for the corresponding sentence.

A core technology of GSEs is Retrieval-Augmented Generation (RAG), which retrieves relevant content from outside of LLMs and incorporates it into LLM answer generation~\cite{internet-augmented-language-models-fewshot-arxiv22, internet-augmented-dialogue-generation-arxiv21, rag-kdd24}.
Information retrieval systems that use RAG are called \emph{RAG systems}~\cite{mirage-acl24, lála2023paperqa, almanac-nejmai-2024}.

However, several studies of \emph{RAG systems} implicitly focus on systems that retrieve from fixed, curated, and closed external databases defined by developers~\cite{mirage-acl24, lála2023paperqa, almanac-nejmai-2024}.
While GSEs similarly use the web as an external database, the nature of this database differs from that of these RAG systems.
The web targeted by GSEs is open, allowing anyone to freely publish information, and is dynamic, with information being updated as needed.
Therefore, unlike these RAG systems, GSEs use search engines (SE) for retrieval from the external database (web).
In our study, we focus on the differences in external database characteristics between GSEs and the RAG systems.
